\documentclass[11pt]{article}
\usepackage[margin=1.15in]{geometry}
\usepackage{amsmath,amssymb,amsthm}
\usepackage{xcolor}
\usepackage[colorlinks=true,linkcolor=blue!55!black,citecolor=blue!55!black,urlcolor=blue!55!black]{hyperref}

\newtheorem{theorem}{Theorem}
\newtheorem{lemma}{Lemma}
\newtheorem{corollary}{Corollary}
\theoremstyle{definition}
\newtheorem{definition}{Definition}
\theoremstyle{remark}
\newtheorem{remark}{Remark}

\DeclareMathOperator{\disc}{disc}
\DeclareMathOperator{\Res}{Res}
\newcommand{\eps}{\varepsilon}

\title{\bfseries A Discriminant Lower Bound for Divisibility-Covering\\
\large Difference Sets, with an Application to the Covering-Pair Exponents}
\author{Eric Schultz \and Claude (Anthropic)}
\date{July 2026}

\begin{document}
\maketitle

\begin{abstract}
Let $S$ be a finite set of integers with all elements of absolute value at most
$M$. We say $S$ \emph{divisibility-covers} a set $I$ of positive integers if
every $i\in I$ divides some nonzero difference $s-s'$ with $s,s'\in S$. We prove,
by an elementary discriminant argument, that covering the primes in a dyadic
interval $(n/2,n]$ forces
$|S|^2\log M \gtrsim n/2$. In the exponent scaling $|S|=n^{\beta}$,
$M=\exp(n^{\alpha})$, this reads $2\beta+\alpha\ge 1$. The proof identifies the
covering condition with divisibility of the Vandermonde discriminant of the
polynomial $\prod_{s\in S}(x-s)$ by the primorial of the interval, and the bound
is the magnitude comparison $|\!\disc| \ge \text{primorial}$. We include a remark
delimiting what the method does and does not give: it is essentially tight as a
\emph{magnitude} bound and cannot by itself yield any stronger exponent
inequality.
\end{abstract}

\section{Statement}

For a finite set $S\subset\mathbb Z$ write $\Delta(S)=\{\,|s-s'| : s,s'\in S,\ s\ne s'\,\}$
for its set of positive differences.

\begin{definition}
Let $I\subset\mathbb Z_{>0}$. We say $S$ \emph{divisibility-covers} $I$ if for
every $i\in I$ there exists $d\in\Delta(S)$ with $i\mid d$.
\end{definition}

Throughout, $n\ge 4$ and $M\ge 2$ are real parameters, and
\[
  P(n) = \{\,p \text{ prime} : n/2 < p \le n\,\},\qquad
  Q(n) = \prod_{p\in P(n)} p .
\]

\begin{theorem}\label{thm:main}
Let $S\subset\mathbb Z$ be finite with $|s|\le M$ for all $s\in S$, and suppose
$S$ divisibility-covers $P(n)$. Then
\[
  |S|\,(|S|-1)\,\log(2M)\ \ge\ \log Q(n)\ =\ \theta(n)-\theta(n/2),
\]
where $\theta$ is the first Chebyshev function. In particular, since
$\theta(n)-\theta(n/2)=\tfrac{n}{2}+o(n)$,
\[
  |S|^2\ \ge\ \frac{n}{2\log(2M)}\,\bigl(1+o(1)\bigr).
\]
\end{theorem}

\begin{corollary}\label{cor:exponents}
Suppose a family of covering sets satisfies $|S|=n^{\beta+o(1)}$ and
$M=\exp\!\bigl(n^{\alpha+o(1)}\bigr)$ with $0\le\alpha<1$. Then
\[
  2\beta+\alpha\ \ge\ 1 .
\]
\end{corollary}

\section{Proof}

Let $S=\{s_1,\dots,s_K\}$ with $K=|S|$ and form the monic polynomial
\[
  f(x)=\prod_{k=1}^{K}(x-s_k)\ \in\ \mathbb Z[x].
\]
By the product-and-sum rule, $f'(x)=\sum_{k}\prod_{j\ne k}(x-s_j)$, and the
discriminant of $f$ satisfies the classical identities
\begin{equation}\label{eq:disc}
  \disc(f)\ =\ (-1)^{K(K-1)/2}\Res(f,f')\ =\ \prod_{i<j}(s_i-s_j)^2 .
\end{equation}

\begin{lemma}\label{lem:equiv}
For a prime $p$, we have $p\mid \disc(f)$ if and only if some nonzero difference
$s_i-s_j$ is divisible by $p$; equivalently, iff $f$ has a repeated root modulo
$p$.
\end{lemma}

\begin{proof}
By \eqref{eq:disc}, $p\mid\disc(f)$ iff $p\mid\prod_{i<j}(s_i-s_j)^2$, iff
$p\mid s_i-s_j$ for some $i<j$ (as $p$ is prime). The same condition says two
roots of $f$ coincide modulo $p$, i.e.\ $f$ has a repeated root mod $p$.
\end{proof}

Now suppose $S$ divisibility-covers $P(n)$. Fix $p\in P(n)$. Since $p>n/2\ge 2$
and $p$ divides some $d=|s_i-s_j|\in\Delta(S)$, Lemma~\ref{lem:equiv} gives
$p\mid\disc(f)$. As the primes in $P(n)$ are distinct, their product divides
$\disc(f)$:
\begin{equation}\label{eq:Qdiv}
  Q(n)=\prod_{p\in P(n)}p\ \ \Big|\ \ \disc(f).
\end{equation}

We now compare magnitudes. First, $\disc(f)\ne 0$: the $s_k$ are distinct
integers, so every factor $s_i-s_j$ in \eqref{eq:disc} is nonzero, hence
$\disc(f)$ is a nonzero integer. From \eqref{eq:Qdiv} and $\disc(f)\ne 0$,
\begin{equation}\label{eq:magnitude}
  |\disc(f)|\ \ge\ Q(n).
\end{equation}
Second, each $|s_i-s_j|\le |s_i|+|s_j|\le 2M$, so by \eqref{eq:disc},
\begin{equation}\label{eq:discub}
  \log|\disc(f)|\ =\ 2\sum_{i<j}\log|s_i-s_j|\ \le\ 2\binom{K}{2}\log(2M)\ =\ K(K-1)\log(2M).
\end{equation}
Combining \eqref{eq:magnitude} and \eqref{eq:discub},
\[
  K(K-1)\log(2M)\ \ge\ \log Q(n)\ =\ \sum_{p\in P(n)}\log p\ =\ \theta(n)-\theta(n/2),
\]
which is the first display of the theorem. By the prime number theorem in the
form $\theta(x)=x+o(x)$ we have $\theta(n)-\theta(n/2)=\tfrac n2+o(n)$, giving
$K^2\ge \dfrac{n}{2\log(2M)}(1+o(1))$.

For Corollary~\ref{cor:exponents}, substitute $K=n^{\beta+o(1)}$ and
$\log(2M)=n^{\alpha+o(1)}$: then $n^{2\beta+o(1)}\ge n^{1-\alpha+o(1)}$, i.e.
$2\beta+\alpha\ge 1$. \qed

\section{Scope of the method}

\begin{remark}[The bound is tight \emph{as a magnitude bound}]
Inequality \eqref{eq:discub} is sharp to leading order: for $K$ points spread in
$[0,M]$, the maximum of $\sum_{i<j}\log|s_i-s_j|$ is $\tfrac{K^2}{2}\log M
-O(K^2)$ (attained near Fekete/Chebyshev configurations). Hence no choice of the
positions of the $s_k$ improves the exponent in Theorem~\ref{thm:main}: the
additive structure of a difference set affects only the $O(K^2)$ lower-order
term, not the leading $\tfrac{K^2}{2}\log M$. Consequently the discriminant
magnitude comparison \emph{cannot} yield any exponent bound stronger than
$2\beta+\alpha\ge1$.
\end{remark}

\begin{remark}[What is \emph{not} proved]
Theorem~\ref{thm:main} is consistent with, but does not establish, the stronger
inequality $\alpha+\beta\ge 1$ (which would require $\beta\ge 1-\alpha$, roughly
twice the exponent proved here). Whether $2\beta+\alpha\ge1$ is attained --
equivalently, whether divisibility-covering sets exist with
$|S|=n^{(1-\alpha)/2+o(1)}$ -- is left open. Establishing either a matching
construction or a strictly stronger lower bound appears to need input beyond the
magnitude of the discriminant; heuristic considerations (a sum--product
antagonism between additive economy and the multiplicative richness of a
difference set) suggest the honest frontier lies there.
\end{remark}

\begin{remark}[Covering the full interval]
If $S$ divisibility-covers all of $(n/2,n]\cap\mathbb Z$ (not merely the primes),
the hypothesis of Theorem~\ref{thm:main} holds a fortiori, so the same lower
bound applies.
\end{remark}

\end{document}
